\backhead{FURTHER EXERCISES}{khred}
\begin{exlist}
\item Five yearly drainage values from Daisy are 150, 210, 95, 180, 130 mm and from a simple model 165, 190, 120, 170, 150 mm. Find the RMSE, the bias and the Nash--Sutcliffe efficiency.\exans{18.7 mm; +6.0 mm; 0.78}
\item A water balance gives inputs of 755 mm and outputs of 742 mm with a storage rise of 4 mm. Find the residual and the relative closure error. Is it acceptable if the tolerance is 1 per cent of the inputs?\exans{9 mm; 1.2 per cent; no}
\item Of 10\,190 runs, 149 failed. What fraction is valid? If failures were concentrated in one experiment of 740 runs with 120 failures, what fraction of that experiment failed, and why does it matter for the conclusions?\exans{0.985; 0.162; the failures are not random}
\item A law has an in-sample RMSE of 0.03 and a leave-one-soil-out RMSE of 0.12. A second law has 0.05 and 0.06. Which should be used, and why?\exans{the second; the first is overfitted}
\item Split the mean squared error of a model with $\overline O=300$, $\overline P=285$, $s_O=60$, $s_P=50$ and $r=0.9$ into bias, scale and scatter.\exans{225, 16, 684; MSE 925}
\end{exlist}
\backhead{ESSAY QUESTIONS}{khred}
{\small\itshape Answer each question in about one page, with equations, a sketch and at least one number from this chapter.}\par\vspace{3pt}
\begin{exlist}
\item What does it mean to audit a model? Describe the steps of the audit of the water balance in this book and what each step can and cannot reveal.
\item Discuss the unattributed term of the water balance: how it arises, how it was detected and why it matters for the interpretation of evaporation savings.
\item Compare in-sample fit with leave-one-out prediction. Why is the second the stronger test of a law, and when can even it mislead?
\item Explain why failed runs should be kept and analysed rather than deleted. Give examples of what failures can teach about a model.
\item Discuss the difference between verification and validation of a model and give one example of each from SWAFL.
\item A compact law reproduces evapotranspiration of an ensemble with high skill. Discuss what such a law is useful for and where it must not be used.
\end{exlist}
\dykband{EXERCISES IN SWAFL}
{\small Open SWAFL (desktop or online) with the benchmark project. The chapter of the \href{https://github.com/halilunaziru73-creator/Coupled-Water-Fluxes-Pore-Scale-Dynamics-Precision-Irrigation-Mulching-Feedback-Loops/raw/main/Software/v1.1.0/SWAFL_Desktop/docs/manual_v1.1.0/SWAFL_Manual_v1.1.0.pdf}{SWAFL manual} that explains each task is given in brackets; record your values in a table and compare them with what is expected (\emph{Check}).}\par\vspace{3pt}
\begin{exlist}
\item {\color{khblue}[Manual, Chapter 14]} Run the field balance of the loam and open the \emph{Messages} tab of the Dashboard. Note every warning about the closure of the balance and explain it.\par\noindent{\small\itshape Check: no warning when the balance closes within 1 mm.}
\item {\color{khblue}[Manual, Chapter 30]} In the Daisy results open \emph{Processes, plots, balance and tests} for one experiment and read the mass balance of water and nitrogen.\par\noindent{\small\itshape Check: the residuals are small compared with the fluxes.}
\item {\color{khblue}[Manual, Chapter 30]} Open \emph{Why a run gave no results} for an experiment with a failed run and find the line of the set-up that caused it.\par\noindent{\small\itshape Check: the cause is located at a line of the set-up file.}
\item {\color{khblue}[Manual, Chapter 22]} Open \emph{Full water balance} in the seasonal data and check the closure of three seasons by hand.\par\noindent{\small\itshape Check: each season closes within rounding.}
\end{exlist}
